\section{Positive matrices at the finite spectral nodes}\label{sa:finite}
We use the cofactors $R_j$ and interval polynomials $K_{i,j}$ of
\eqref{lc:R-recurrence}--\eqref{sa:path}. Throughout the remaining
sections, $u(\delta)=-4\sin^2(\pi\delta)$ denotes the spectral coordinate;
it is distinct from the protected vertex $u$ of the triangulation.
We also write $P_{\mathrm C}=\mathcal P_{\mathrm C}=(-1)^dH_{V,\mathrm C}$.
The central factorizations \eqref{lc:even-factor}--\eqref{lc:odd-factor}
retain the odd skipped-row term, including its derivative at $u=0$.

We give the finite positive representation needed for phase estimates.
Let $S_j=\prod_{r=1}^j s_r$, $S_0=1$, and, for $0\le a\le J$,
let $H_\pm^{(a)}$ be the Hermitian tridiagonal matrices with
\begin{align}
 (H_\pm)_{n,n}&=\frac1{\cos(n\theta)\cos((n+1)\theta)},
                              &&0\le n<a,\notag\\
 (H_\pm)_{n-1,n}&=\sqrt{b_n}\{\sin(n\theta)\pm\ii\cos(n\theta)\},
                              &&1\le n<a,\label{sa:H}\\[-6pt]
 b_n&=\frac{\sin^2(n\theta)}{\cos^2(n\theta)
                  \sin((2n-1)\theta)\sin((2n+1)\theta)}.\notag
\end{align}
Empty matrices have determinant one. Put
\begin{equation}\label{sa:D}
 \mathcal D_a(x,b,y)=\det[bI+xH_-^{(a)}/2+yH_+^{(a)}/2].
\end{equation}

\begin{lemma}[Exact node adapter]\label{sa:node}
For $u_a=-a_a$ and $a\le j\le J$,
\begin{equation}\label{sa:node-id}
 R_a(b,u_a)=S_{2a}\mathcal D_a(x,b,y),\qquad
 R_j(b,u_a)=K_{a,j}(u_a)R_a(b,u_a).
\end{equation}
Both $H_\pm^{(a)}$ are positive definite. The polynomials
$\mathcal D_a$ and $K_{a,j}(u_a)$ have nonnegative face coefficients.
\end{lemma}
\begin{proof}
At $u_a$, the cofactor column on rows $0,\ldots,a$ is a null vector
of the original triangular matrix. If its coefficients are $r_i$, its
row equations are
\[
 (a_i-a_a)r_i=s_{i+1}^2A r_{i+1}+s_{i+2}s_{i+1}P r_{i+2}.
\]
The leading coefficient is
$\prod_{r<a}(a_r-a_a)=S_{2a}$, by
$a_r-a_a=4\sin((a+r+1)\theta)\sin((a-r)\theta)$.
Define $\partial_q b^j=s_j b^{j-1}$. Since
$a_i-a_a=s_as_{a+1}-s_is_{i+1}$, the normalized monic polynomial
is the unique degree-$a$ eigenpolynomial of
\[
 \mathcal L=\partial_q\bigl((b^2+Ab+P)\partial_q\bigr),
                     \qquad\lambda_a=s_as_{a+1}.
\]
The eigenvalues through $J$ are distinct: their successive differences
are $4\sin(2a\theta)\sin\theta>0$.

For completeness its monic eigenpolynomials $p_a$ obey
\begin{equation}\label{sa:Jacobi}
 p_a=\left(b+\frac{A}{2\cos((a-1)\theta)\cos(a\theta)}\right)p_{a-1}
       -\frac{b_{a-1}}4\{A^2-4P\cos^2((a-1)\theta)\}p_{a-2}.
\end{equation}
Here $p_0=1,p_1=b+A/(2\cos\theta)$.
One elementary derivation also checks that no lower-degree terms were
dropped. Normalize $[j]=s_j/s_1$ and define moments by
\[
 \mu_0=1,\quad\mu_{-1}=0,\qquad
 [r+1]\mu_r+A[r]\mu_{r-1}+P[r-1]\mu_{r-2}=0
                              \quad(1\le r\le2J-1).
\]
The sine addition formulas make $\mathcal L$ symmetric for
$\langle b^i,b^j\rangle=\mu_{i+j}$. Distinct eigenvalues therefore
make the $p_a$ orthogonal on the available degrees. Multiplication
by $b$ is symmetric as well, so $bp_{a-1}-p_a$ has only the two
previous components. Comparison of the two highest coefficients,
using the displayed row equations, gives \eqref{sa:Jacobi}.
Their norm recurrence has its nonzero second coefficient for generic
$(A,P)$, starting from norm one. This justifies the induction before
polynomial continuation to all $A,P$; no positive moment assumption
is used. All denominators have index less than $N$.
The tridiagonal determinant \eqref{sa:D} has precisely
\eqref{sa:Jacobi}, proving the first identity.

The Cholesky pivots of either $H$ are
\[
 p_{n-1}^{H}=\frac{\sin(n\theta)}
                     {\cos(n\theta)\sin((2n-1)\theta)}>0,
                         \qquad1\le n\le a.
\]
Writing $d_n=(H_\pm)_{nn}$ and $h_n=(H_\pm)_{n-1,n}$, substitution
in $d_n-|h_n|^2/p_{n-1}^{H}=p_n^{H}$ verifies the
factorization, including the last physical row. Thus both matrices
are positive. Cauchy--Binet applied to their Cholesky factors proves
coefficient positivity of \eqref{sa:D}.
The second identity and coefficient positivity of $K_{a,j}(u_a)$
were proved in Lemma~\ref{lc:nodal-positive}: after row $a$ the forcing
vanishes, and the solutions of the normalized recurrence
\eqref{lc:tail-node} have nonnegative face coefficients despite its
signed second term.
\end{proof}

At the maximal physical fold the scalar can be read without any
interpolation. For $\delta_*=0$ at odd levels and $\delta_*=1/(2N)$
at even levels, \eqref{sa:even}--\eqref{sa:node-id} give
\begin{equation}\label{sa:anchor}
 P_{\mathrm C}(\delta_*)=N^2\beta_J^2
          \mathcal D_J(x,b,y)\mathcal D_J(x,c,y)(x+y)^{d-2J},
                       \qquad \beta_J=S_{2J}/S_J^2.
\end{equation}
Indeed $S_{N-1}=N$, $S_J^2=N$ at odd levels; at even levels
$S_{N-2}=N/(2\sin\theta)$, $S_J^2=N/2$, and $s_m=2$.
Thus both the even forced factor and the odd simple node are literal.

